\documentclass[11pt]{article}
\usepackage[a4paper,margin=25mm]{geometry}
\usepackage{amsmath,amssymb,amsthm,hyperref}
\newtheorem{theorem}{Theorem}
\newtheorem{lemma}{Lemma}
\newcommand{\PP}{\mathcal P}
\title{Polynomial conditioning of a fixed number of cumulative moments}
\author{Proof independently reviewed by root and independent\_directions}
\date{30 September 2026}
\begin{document}\maketitle
A single Legendre perturbation places all its cumulative moments in
the polynomial module spanned by $1,H_d,H_d'$, so it cannot condition
arbitrarily many independent moment features. A fixed number of
separated perturbation bands removes that obstruction.

\begin{theorem}\label{main}
Fix $s\ge1$, let $r\ge0$, and put $R=r+s+1$. Choose positive integers
\[
 d_i>2\left(R+\sum_{j<i}(d_j+1)\right)+1,
 \qquad 1\le i\le s,
\]
and put $D=R+\sum_i(d_i+1)$. For $0<\eta\le1/(2s)$ define the
strictly positive density and its cumulative moments
\[
 f(t)=1-\eta\sum_{i=1}^sP_{d_i}(2t-1),\qquad
 M_l(t)=\int_0^t u^l f(u)du\quad(0\le l\le s).
\]
There are constants $C_s,A_s$ such that, for arbitrary
$p_*,p_0,\ldots,p_s\in\PP_r$,
\[
 \left(\|p_*\|_2^2+\sum_{l=0}^s\|p_l\|_2^2\right)^{1/2}
 \le C_s(D+1)^{A_s}\eta^{-2s}
       \left\|p_*+\sum_{l=0}^sp_lM_l\right\|_2.
\]
All norms are for Lebesgue measure on $[0,1]$. Minimal degree choices
have $D=O_s(r+1)$. The same bound, up to fixed constants, holds in
any measure whose density is bounded above and below by fixed positive
constants. It also holds empirically on a real quadrature exact to
at least twice the degree of the displayed combination.
\end{theorem}
\begin{proof}
Write $w=t(1-t)$, $H_i(t)=\int_0^tP_{d_i}(2u-1)du$, and
$\lambda_i=d_i(d_i+1)$. For $l\ge1$, integration by parts gives
\[
 M_l=\frac{t^{l+1}}{l+1}
 -\eta\sum_i\left[t^lH_i-l\int_0^t u^{l-1}H_i(u)du\right].
\]
For $l=0$, $M_0=t-\eta\sum_iH_i$.
The exact primitive lemma in \path{ordered_prefix_legendre_tools.tex}
uses $\mathcal L_iC=(wC)''+\lambda_iC$ and gives
\[
 \int_0^t u^{l-1}H_i(u)du=A_{il}(t)H_i+B_{il}(t)H_i',
 \qquad B_{il}=-w\mathcal L_i^{-1}t^{l-1}.
\]
Here $\deg A_{il}\le l$ and $\deg B_{il}\le l+1$; their coefficients
and norms are bounded by a fixed power of $D$ for fixed $s$.
Consequently the combination $F=p_*+\sum_lp_lM_l$ has the form
\[
 F=P+\eta\sum_i(H_i A_i+H_i'B_i),\qquad
 \deg P,\deg A_i,\deg B_i\le R,
\]
where, crucially,
\begin{equation}\label{B}
 B_i=-w\sum_{l=1}^s l p_l\mathcal L_i^{-1}t^{l-1}.
\end{equation}
The trinary-band conditioning theorem from the same note yields
\[
 \|P\|_2+\sum_i(\|A_i\|_2+\|B_i\|_2)
 \le C_s(D+1)^{6s}\eta^{-s}\|F\|_2.
\]

We show that the $B_i$ recover $p_1,\ldots,p_s$ polynomially.
Let $J_0,\ldots,J_{s-1}$ be the shifted Jacobi polynomials for weight
$w$, with any fixed rational normalization. They satisfy
\[
 \mathcal L_iJ_j=(\lambda_i-\mu_j)J_j,
 \qquad\mu_j=(j+1)(j+2).
\]
Expand $t^{l-1}=\sum_{j=0}^{l-1}a_{jl}J_j$ and set
$c_j=\sum_{l=j+1}^s l a_{jl}p_l$. This is a fixed invertible
triangular transformation from $(p_1,\ldots,p_s)$ to $(c_0,\ldots,c_{s-1})$.
Equation~\eqref{B} becomes
\[
 -B_i=\sum_{j=0}^{s-1}\frac{wJ_j c_j}{\lambda_i-\mu_j}.
\]
The scalar Cauchy matrix $C_{ij}=1/(\lambda_i-\mu_j)$ is invertible.
Its determinant has absolute value
\[
 \frac{\prod_{i<i'}|\lambda_{i'}-\lambda_i|
             \prod_{j<j'}|\mu_{j'}-\mu_j|}
      {\prod_{i,j}|\lambda_i-\mu_j|}.
\]
All differences are nonzero integers, and each denominator factor
is at most $C_s(D+1)^2$. Thus the determinant is bounded below by
an inverse polynomial in $D$, while all entries and cofactors are
bounded by constants depending only on $s$. The inverse matrix has
polynomial norm. It follows that each $\|wJ_j c_j\|_2$ is at most
$\operatorname{poly}_s(D)\eta^{-s}\|F\|_2$.

Multiplication by any fixed nonzero polynomial $Q$ of degree at most
$s+1$ has a polynomial lower bound on $\PP_r$:
\[
 \|Qp\|_2\ge c_Q(r+1)^{-2(s+1)}\|p\|_2.
\]
To see this, remove intervals of length $c_Q(r+1)^{-2}$ around its
finitely many real roots in $[0,1]$. The polynomial evaluation bound
retains at least half the squared norm of $p$ on the remainder, and
factorization of $Q$ gives the stated pointwise lower bound there.
Apply this to the fixed polynomials $wJ_j$. We recover all $c_j$,
and therefore all $p_l$, $l\ge1$, with polynomial loss.

Since $f$ is a probability density between $1/2$ and $3/2$,
$0\le M_l\le1$. Subtracting $\sum_{l=1}^sp_lM_l$ from $F$
therefore leaves $p_*+p_0M_0$ with polynomially bounded norm.
Its trinary representation has constant coefficient $p_*+tp_0$
and each $H_i$ coefficient $-p_0$. One more application of the
same conditioning theorem bounds $p_0$ and then $p_*$, at the cost
of another $\eta^{-s}$. This proves the displayed result. The
weighted and exact-quadrature statements follow immediately by norm
comparison and exactness for the squared polynomial.
\end{proof}

\paragraph{Moment compatibility against a polynomial density.}
Suppose an additional density $v$ has degree $e$, and take the first
band degree larger than $e+m+s+1$. For $0\le a\le m$ and $0\le l\le s$,
Fubini and Legendre orthogonality give
\[
 \int_0^1 t^a M_l(t)v(t)dt
 =\frac1{l+1}\int_0^1t^{a+l+1}v(t)dt.
\]
Indeed the perturbation integral is against $P_{d_i}(2u-1)$ times
$u^l\int_u^1t^av(t)dt$, whose degree is at most $e+a+l+1<d_i$.
This provides a polynomially conditioned real baseline for rational
cumulative-moment ranks. Their constant targets depend only on the
first $s+1$ moments of $v$.

\paragraph{Arithmetic limitation for iterated constructions.}
The fixed-$s$ box-smoothing theorem additionally preserving its first
$s$ conditional moments needs targets with $a\le s$ in the preceding
identity, hence input moments through degree $2s+1$. Furthermore,
polynomial denominator bounds for separate local constructions do not
imply a polynomial \emph{common} denominator across polynomially many
cells. Synchronizing those denominators is a separate requirement;
no construction with arbitrarily many fixed small dice is inferred here.
\end{document}
